% Generated draft: requires mathematical and source review.
\documentclass[11pt]{article}
\usepackage[margin=25mm]{geometry}
\usepackage{fontspec}
\setmainfont{Latin Modern Roman}
\usepackage{amsmath,amssymb,mathtools,unicode-math}
\setmathfont{Latin Modern Math}
\usepackage{enumitem,needspace,xurl,hyperref}
\setlength{\parindent}{0pt}
\setlength{\parskip}{5pt}
\setlength{\emergencystretch}{3em}
\newcommand{\N}{\mathbb{N}}
\newcommand{\Z}{\mathbb{Z}}
\newcommand{\Q}{\mathbb{Q}}
\newcommand{\R}{\mathbb{R}}
\newcommand{\C}{\mathbb{C}}
\newcommand{\F}{\mathbb{F}}
\newcommand{\A}{\mathbb{A}}
\newcommand{\PP}{\mathbb{P}}
\newcommand{\E}{\mathbb{E}}
\newcommand{\eps}{\varepsilon}
\newcommand{\dd}{\,\mathrm{d}}
\begin{document}
\section*{UnsolvedMath ID: 30003794}
\textbf{Source problem code:} catalogue entry
\subsection{Definitions and mathematical statement}
Let \(K\) be a field and let \(G\) be an absolutely almost simple algebraic \(K\)-group that is either simply connected or adjoint. Here absolutely almost simple means that, over an algebraic closure of \(K\), the group is nontrivial, connected, semisimple, and has no proper nontrivial connected normal algebraic subgroup. Simply connected is understood in the sense of semisimple algebraic groups; adjoint means that the center is trivial.

An \emph{inner \(K\)-form} of \(G\) is an algebraic \(K\)-group \(H\) admitting an isomorphism \(f:H_{K_s}\to G_{K_s}\), where \(K_s\) is a separable closure of \(K\), such that the descent cocycle \(f\,{}^{\sigma}\!f^{-1}\) takes values in \(\operatorname{Inn}(G)(K_s)\) for every \(\sigma\in\operatorname{Gal}(K_s/K)\). The groups \(H\) are counted up to \(K\)-isomorphism.

For a nontrivial discrete valuation \(v\) of \(K\), write \(K_v\) for the completion and \(\mathcal O_v\) for its valuation ring. Using the standard completed local convention, \(H\) has \emph{good reduction at \(v\)} if there exists a reductive group scheme \(\mathcal H\) over \(\mathcal O_v\) with
\[
\mathcal H\times_{\mathcal O_v}K_v\simeq H\times_K K_v.
\]
A reductive group scheme here is a smooth affine group scheme of finite presentation whose geometric fibers are connected reductive algebraic groups. The supplied record does not specify its local-model convention; the completed local convention is made explicit here for review against the cited source.

\textbf{Problem.} Determine under what hypotheses on \(K\) and \(G\) the field \(K\) admits a natural set \(V\) of discrete valuations such that
\[
\text{for every finite }S\subseteq V,\qquad
\#\left\{
[H]_K\;\middle|\;
\begin{gathered}
H\text{ is an inner }K\text{-form of }G,\\
H\text{ has good reduction at every }v\in V\setminus S
\end{gathered}
\right\}<\infty.
\]
The set \(V\) is chosen before the quantification over \(S\). No restriction on the characteristic of \(K\), or assumption that \(K\) is finitely generated, is stated in the problem. The word ``natural'' is not defined in the supplied record, so the admissible choices of \(V\), and consequently the precise classification target, remain unresolved. In particular, this is a request to determine circumstances in which the property holds, rather than an assertion that it holds for every field and group.
\subsection{Short English statement}
Determine when a field \(K\) and an absolutely almost simple, simply connected or adjoint \(K\)-group \(G\) admit a natural collection of discrete valuations for which only finitely many inner forms of \(G\), up to \(K\)-isomorphism, have good reduction outside any prescribed finite subset. The source leaves ``natural'' unspecified. Its supplied label is partially solved; this draft does not independently verify that assessment.
\subsection{Sources}
\begin{itemize}
\item Supplied local registry record, permanent UnsolvedMath ID \(30003794\), source problem code \texttt{catalogue entry}, title \emph{Finiteness of Algebraic Forms with Good Reduction}, in the Oberwolfach Reports Open Problems collection. No individual proposer is supplied; the proposed year is recorded as 2018.
\item Supplied report citation: \emph{Quadratic Forms and Related Structures over Fields} (2019), Oberwolfach Reports, pp.~1235--1291, DOI \texttt{10.4171/owr/2018/21}. \url{https://doi.org/10.4171/owr/2018/21}. The background attributes the problem to workshop open-problem discussions documented in Oberwolfach Reports. The citation year 2019 and the record's 2018 year indicators are preserved without reconciliation.
\item Supplied literature source, labeled ``spinor groups with good reduction'': \url{https://arxiv.org/abs/1707.08062}.
\item Supplied literature source, labeled ``good-reduction survey'': \url{https://arxiv.org/abs/2005.05484}.
\item Supplied literature source, labeled ``workshop report'': \url{https://ems.press/content/serial-article-files/46742}.
\item Literature assessment dated exactly \texttt{2026-08-21}: the supplied curation labels the problem partially solved and reports finiteness results in important cases, including inner type \(A\), spinor groups over two-dimensional global fields, and related types, while describing the arbitrary-field and general finitely-generated-field problem as open. This is supplied dated curation, not independently verified. The record expressly requests independent verification of time-sensitive status and the cited claims; the exact hypotheses of those reported results are not supplied.
\item The registry marks the statement \texttt{corrected\_verified} and reports a manual semantic, typographic, and LaTeX review. These are supplied review annotations, not independent mathematical verification by this draft.
\end{itemize}
\end{document}
